Redes varejistas precisam decidir continuamente quando e quanto repor estoque em cada loja. Essa decisão se torna mais difícil quando os produtos não apresentam o mesmo padrão de venda: alguns têm demanda regular, enquanto outros passam vários períodos sem venda e concentram vendas em poucos ciclos. Nesses casos, aplicar uma única regra de reposição para toda a rede pode gerar excesso de estoque em alguns produtos, falta de estoque em outros e maior instabilidade nos pedidos enviados ao armazém.
Esta dissertação propõe um \textit{framework} metodológico de apoio à decisão para recomendar políticas de reposição de acordo com o comportamento observado em cada série loja-produto. A hipótese central é que a política mais adequada depende do regime operacional da demanda, e que essa relação pode ser aprendida a partir de características históricas da série. Para isso, a proposta utiliza como entrada o histórico preparado, a caracterização operacional das séries e sinais auxiliares de previsão de demanda. A partir dessas informações, avalia políticas candidatas em simulação e aprende a associar perfis de demanda a políticas com melhor desempenho esperado. Esse \textit{framework} é denominado \textit{Adaptive Inventory Policy Engine} (AIPE), e seu componente de seleção é o \textit{Policy Selection Engine} (PSE).
A proposta é avaliada com dados reais de uma rede varejista brasileira, cuja base contém aproximadamente 48 milhões de registros transacionais e 12 mil produtos. Nos subconjuntos experimentais já analisados, os resultados indicam redução de custo total de inventário de até 48\% em relação a políticas de referência, sem degradação proporcional do nível de serviço. A avaliação compara políticas clássicas, meta-heurísticas e agentes de aprendizado por reforço.
% Os resultados preliminares indicam que, no conjunto de políticas e séries avaliadas, não há uma política universalmente dominante. Em vez disso, mesmo dentro do regime de demanda altamente irregular (\textit{Lumpy}) analisado nas Fases~1 e~2, diferentes perfis operacionais favorecem políticas distintas; a generalização para os demais regimes de demanda é objetivo da Fase~3. 
A principal contribuição da dissertação é avaliar a viabilidade de selecionar políticas de reposição de forma contextual, com base em evidência empírica, para apoiar decisões de estoque em redes varejistas heterogêneas.

\vspace{2em}
\textbf{Palavras-chave}: Otimização de Inventário; Demanda Intermitente; Seleção Contextual de Políticas; Cadeia de Suprimentos.
